\documentclass[11pt]{article}
\usepackage[utf8]{inputenc}
\usepackage{amsmath, amssymb}
\usepackage{geometry}
\usepackage{listings}
\usepackage{xcolor}
\usepackage{hyperref}
\usepackage{tabularx}
\usepackage{booktabs}

\geometry{a4paper, margin=1in}

\definecolor{codegray}{rgb}{0.5,0.5,0.5}
\definecolor{codepurple}{rgb}{0.58,0,0.82}
\definecolor{backcolour}{rgb}{0.95,0.95,0.92}

\lstset{
    language=Python,
    backgroundcolor=\color{backcolour},
    commentstyle=\color{codegray},
    keywordstyle=\color{magenta},
    stringstyle=\color{codepurple},
    basicstyle=\ttfamily\small,
    breaklines=true,
    showstringspaces=false
}

\title{Optional Problem: Inverse Quantum Tomography \\ \large Differentiable Physics \& Gradient-Based Inversion}
\author{VB}
\date{\today}

\begin{document}

\maketitle

\section*{Prerequisites}

\begin{itemize}
    \item Completed QM Problem 5 (split-operator method)
    \item Basic familiarity with gradient-based optimization
\end{itemize}

\section*{PyTorch Quick Reference}

\begin{center}
\renewcommand{\arraystretch}{1.3}
\begin{tabularx}{\textwidth}{@{} >{\ttfamily\small}l X @{}}
\toprule
\textcolor{blue}{PyTorch} & \textcolor{blue}{NumPy equivalent / Notes} \\ \midrule
import torch & import numpy as np \\
torch.tensor(arr) & Creates a tensor from a NumPy array \\
t.numpy() & Converts tensor back to NumPy \\
torch.fft.fft(t) / ifft(t) & Same as numpy.fft.fft / ifft \\
torch.fft.fftfreq(n, d=dx) & Same as numpy.fft.fftfreq \\
torch.exp, torch.sin, \ldots & Element-wise, same as NumPy \\
t.requires\_grad\_(True) & Enable gradient tracking on tensor \\
loss.backward() & Compute $\nabla_\theta \mathcal{L}$ via backprop \\
optimizer.step() & $\theta \leftarrow \theta - \alpha \nabla_\theta \mathcal{L}$ \\
torch.no\_grad() & Context manager: disable autograd \\
\bottomrule
\end{tabularx}
\end{center}

\textbf{Install:} \lstinline|pip install torch| \quad
\textbf{Docs:} \url{https://pytorch.org/docs/stable/}

\textbf{Key idea:} PyTorch tensors track every operation. Calling \lstinline|loss.backward()| differentiates the loss with respect to all tensors that have \lstinline|requires_grad=True|. This lets us differentiate \emph{through} a physics simulation.


\section*{Problem: Reconstructing an Unknown Potential}

You built a split-operator solver (QM Problem~5). Now imagine you are an experimentalist: you observe a wavepacket scattering off an \textbf{unknown} potential $V_{\text{true}}(x)$, and you want to reconstruct it from measurements.

\subsection*{Mathematical Background}

\textbf{The Inverse Problem.} Given sparse observations $\{|\psi(x_j, t_k)|^2\}$ at various times, find $V(x)$ such that propagating $\psi_0$ through $V$ reproduces the data. The map $V \mapsto \psi(\cdot, t)$ is given by $e^{-i\hat{H}t}$, which depends on $V$ nonlinearly.

\textbf{Teacher Forcing.} Naively unrolling the full simulation and backpropagating causes \emph{cycle skipping}: if $V_\theta$ is slightly wrong, the predicted wave drifts out of phase, and the gradient points the wrong way. The fix: break the trajectory into short segments and \textbf{reset to the true state} at each segment boundary. This bounds phase error and convexifies the loss landscape.


\subsection*{Implementation}

Open \lstinline|implementation_tasks.py|. Implement:

\begin{enumerate}
    \item \textbf{\lstinline|pt_split_operator_step(psi, k_sq, V, dt)|} — the same Trotter step as Problem~5, but in PyTorch (\lstinline|torch.fft.fft/ifft|, \lstinline|torch.exp|). No \lstinline|.numpy()| calls!
    \item \textbf{\lstinline|GridPotential.__init__ / forward|} — learnable $V_\theta(x)$ as an \lstinline|nn.Parameter| smoothed by a Gaussian kernel.
    \item \textbf{\lstinline|run_teacher_forcing(...)|} — for each consecutive pair of snapshots, propagate from the true state and accumulate $\mathcal{L} = \|\psi_\text{pred} - \psi_\text{true}\|^2$.
\end{enumerate}

\subsection*{Test}

\begin{lstlisting}[language=bash]
python levels/level_1_tomography.py
\end{lstlisting}

The learned $V_\theta(x)$ should converge to the hidden $V_\text{true}(x)$ within 100--300 iterations.

\end{document}
